\exercisesection{}

\begin{enumerate}
\def\labelenumi{\arabic{enumi})}
\tightlist
\item
  Let \(\mathbf{E}\) be a set, \(\Phi\) a clan of subsets of \(\mathbf{E}\) , and \(m\) an additive set function defined on \(\Phi\) . Denote by \(\mathcal{E}(\Phi)\) the vector space of real \(\Phi\) -step functions on \(\mathbf{E}\) , and by \(J\) the linear form on \(\mathcal{E}(\Phi)\) such that \(J(\varphi_{\mathbf{A}}) = m(\mathbf{A})\) for all \(A \in \Phi\) .
\end{enumerate}

\begin{enumerate}
\def\labelenumi{\alph{enumi})}
\tightlist
\item
  For every \(\mathbf{A} \in \Phi\) , one sets
\end{enumerate}

\[
m _ {+} (\mathrm{A}) = \sup _ {\mathrm{B} \in \Phi , \mathrm{B} \subset \mathrm{A}} m (\mathrm{B}) \quad \text { and } \quad m _ {-} (\mathrm{A}) = \sup _ {\mathrm{B} \in \Phi , \mathrm{B} \subset \mathrm{A}} (- m (\mathrm{B})).
\]

Show that \(|m|(\mathbf{A}) = m_{+}(\mathbf{A}) + m_{-}(\mathbf{A})\) is the supremum of the set of numbers of the form \(\sum_{i=1}^{p}|m(\mathbf{A}_i)|\) for all the finite partitions \((\mathbf{A}_i)_{1\leqslant i\leqslant p}\) of A into sets belonging to \(\Phi\) .

\begin{enumerate}
\def\labelenumi{\alph{enumi})}
\setcounter{enumi}{1}
\item
  One says that \(m\) is relatively bounded if the set of numbers of the form \(m(\mathbf{B})\) with \(\mathbf{B} \in \Phi\) , \(\mathbf{B} \subset \mathbf{A}\) is bounded for any set \(\mathbf{A} \in \Phi\) . Show that the following conditions are equivalent: \(\alpha)\) \(m\) is relatively bounded; \(\beta)\) the linear form \(J\) on the Riesz space \(\mathcal{E}(\Phi)\) is relatively bounded; \(\gamma)\) \(m\) is the difference of two positive additive set functions on \(\Phi\) ; \(\delta)\) \(|m|(\mathbf{A})\) is finite for all \(\mathbf{A} \in \Phi\) (argue in a circle \(\alpha) \Rightarrow \beta) \Rightarrow \gamma) \Rightarrow \delta) \Rightarrow \alpha)\)).
\item
  Suppose that \(m\) is relatively bounded. Show that the set functions \(|m|\) , \(m_{+}\) and \(m_{-}\) are additive and that \(m = m_{+} - m_{-}\) (make use of Th. 1 of Ch. II, §2, No.~2). Moreover, if \(m'\) and \(m''\) are additive and positive set functions on \(\Phi\) such that \(m = m' - m''\) , then \(m' \geqslant m_{+}\) and \(m'' \geqslant m_{-}\) .
\end{enumerate}

\(d)\) One says that \(m\) is countably additive if \(m(A) = \sum_{n\geqslant 1}m(A_n)\) for every countable

partition \((\mathbf{A}_n)_{n\geqslant 1}\) of a set \(\mathbf{A}\in \Phi\) into sets belonging to \(\Phi\) . Show that this condition signifies that \(\lim_{n\to \infty}m(\mathbf{A}_n) = 0\) for every decreasing sequence \((\mathbf{A}_n)_{n\geqslant 1}\) of elements of \(\Phi\) such that \(\bigcap_{n\geqslant 1}\mathbf{A}_n = \emptyset\) . If \(m\) is countably additive, show that \(|m|\) is countably additive.

\begin{enumerate}
\def\labelenumi{\arabic{enumi})}
\setcounter{enumi}{1}
\tightlist
\item
  Let E be a set, \(\Phi\) a clan of subsets of E, and V a fully lattice-ordered Riesz space. A mapping \(m\) of \(\Phi\) into V is said to be positive if \(m(A) \geqslant 0\) for all \(A \in \Phi\) , that it is additive if \(m(A \cup B) = m(A) + m(B)\) when the sets \(A \in \Phi\) and \(B \in \Phi\) are disjoint, and that it is relatively bounded if the set of elements \(m(B)\) for \(B \in \Phi\) , \(B \subset A\) is bounded above in V for any \(A \in \Phi\) .
\end{enumerate}

\begin{enumerate}
\def\labelenumi{\alph{enumi})}
\tightlist
\item
  Let \(m\) be a relatively bounded mapping of \(\Phi\) into \(V\) . One defines for every \(A\) the elements \(m_{+}(A), m_{-}(A)\) and \(|m|(A)\) of \(V\) as in Exer. 1 a). Show that \(|m|(A)\)
\end{enumerate}

is the supremum of the set of elements of V of the form \(\sum_{i=1}^{p}|m(A_i)|\) for all the finite partitions \((\mathbf{A}_i)_{1\leqslant i\leqslant p}\) of A into sets belonging to \(\Phi\) . From this, deduce that the mapping \(|m|\) of \(\Phi\) into V is additive, then that \(m\) is the difference of the additive positive mappings \(m_{+}\) and \(m_{-}\) of \(\Phi\) into V.

\begin{enumerate}
\def\labelenumi{\alph{enumi})}
\setcounter{enumi}{1}
\tightlist
\item
  Let \(m'\) and \(m''\) be two additive positive mappings of \(\Phi\) into \(V\) and let \(m = m' - m''\) . Show that \(m\) is relatively bounded, and that \(m' \geqslant m_+\) and \(m'' \geqslant m_-\) . c) Generalize Exer. 1 d).
\end{enumerate}

\begin{enumerate}
\def\labelenumi{\arabic{enumi})}
\setcounter{enumi}{2}
\tightlist
\item
  Let \(\mathbf{E}\) be a set and \(\mathfrak{T}\) a tribe of subsets of \(\mathbf{E}\) . Denote by \(\mathcal{M}\) the set of mappings \(f\) of \(\mathbf{E}\) into \(\overline{\mathbf{R}}_+\) such that the set of \(x \in \mathbf{E}\) for which \(f(x) \geqslant c\) belongs to \(\mathfrak{T}\) for every \(c \in \overline{\mathbf{R}}_+\) .
\end{enumerate}

\begin{enumerate}
\def\labelenumi{\alph{enumi})}
\item
  Show that the limit superior and limit inferior of every sequence of elements of \(\mathcal{M}\) belong to \(\mathcal{M}\) (first study sequences that are increasing or decreasing).
\item
  Show that \(\mathfrak{T}\) is the set of subsets of \(\mathbf{E}\) whose characteristic function belongs to \(\mathcal{M}\) .
\item
  Let \(f \in \mathcal{M}\) . For every Borel subset \(B\) of \(\overline{\mathbf{R}}_+\) , the set \(f^{-1}(B)\) belongs to \(\mathfrak{T}\) (the set of subsets \(B\) of \(\overline{\mathbf{R}}_+\) such that \(f^{-1}(B) \in \mathfrak{T}\) is a tribe, containing the intervals \([c, +\infty]\) ).
\item
  Let \(f_1, \ldots, f_n\) be in \(\mathcal{M}\) and let \(\varphi\) be a Borel mapping of \((\overline{\mathbf{R}}_+)^n\) into \(\overline{\mathbf{R}}_+\) ; show that the mapping \(x \mapsto \varphi(f_1(x), \ldots, f_n(x))\) of E into \(\overline{\mathbf{R}}_+\) belongs to \(\mathcal{M}\) . Deduce from this that \(\mathcal{M}\) contains the sum of every series of elements of \(\mathcal{M}\) .
\item
  Show that \(\mathcal{M}\) is the set of limits of increasing sequences of finite positive \(\mathfrak{T}\) -step functions.
\end{enumerate}

\begin{enumerate}
\def\labelenumi{\arabic{enumi})}
\setcounter{enumi}{3}
\tightlist
\item
  Notations and hypotheses are those of the preceding exercise. One calls abstract measure on \((\mathbf{E},\mathfrak{T})\) every mapping \(m\) of \(\mathfrak{T}\) into \(\overline{\mathbf{R}}_+\) having the following property: for every sequence \((\mathrm{A}_n)_{n\in \mathbf{N}}\) of sets belonging to \(\mathfrak{T}\) , pairwise disjoint, one has \(m\big(\bigcup_{n\in \mathbf{N}}\mathrm{A}_n\big) =\)
\end{enumerate}

\(\sum_{n\in \mathbf{N}}m(\mathrm{A}_n)\) . One calls integral on \((\operatorname {E},\mathfrak{T})\) every mapping J of \(\mathcal{M}\) into \(\overline{\mathbf{R}}_+\) satisfying

  the following conditions: \(\alpha)\) \(J(\lambda \cdot f) = \lambda \cdot J(f)\) for \(\lambda \in \overline{\mathbf{R}}_{+}\) and \(f \in \mathcal{M}\) (with the usual convention \(0 \cdot (+\infty) = 0\) ); \(\beta)\) \(J\big(\sum_{n \in \mathbf{N}} f_n\big) = \sum_{n \in \mathbf{N}} J(f_n)\) for every sequence \((f_n)_{n \in \mathbf{N}}\) of elements of \(\mathcal{M}\) .


\begin{enumerate}
\def\labelenumi{\alph{enumi})}
\item
  Let \(J\) be an integral on \((\mathbf{E},\mathfrak{T})\) ; for every subset \(A \in \mathfrak{T}\) , set \(m_{J}(A) = J(\varphi_A)\) ; show that \(m_J\) is an abstract measure and that the mapping \(J \mapsto m_J\) is a bijection of the set of integrals on \((\mathbf{E},\mathfrak{T})\) onto the set of abstract measures on \((\mathbf{E},\mathfrak{T})\) (if \(m\) is an abstract measure, first define \(J(f)\) by linearity for the finite positive \(\mathfrak{T}\) -step functions and use Exer. 3 e)). If \(m\) is an abstract measure, we will denote by \(m^*\) the corresponding integral.
\item
  Let \(m\) be an abstract measure on (E, T). For every positive function \(f\) on E, finite or not, set \(m^{*}(f) = \inf_{g \in \mathcal{M}, g \geqslant f} m^{*}(g)\) . Show that \(m^{*}\) is an encumbrance on E (imitate the proof of Th. 3 of Ch. IV, §1, No.~3). One sets \(m^{*}(A) = m^{*}(\varphi_{A})\) for every subset A of E.
\end{enumerate}

\begin{enumerate}
\def\labelenumi{\arabic{enumi})}
\setcounter{enumi}{4}
\tightlist
\item
  Let \(\mathbf{E}\) be a set, \(\mathfrak{T}\) a tribe of subsets of \(\mathbf{E}\) , \(m\) an abstract measure on \((\mathbf{E},\mathfrak{T})\) , and \(p\geqslant 1\) a finite real number. For every numerical function \(f\) on \(\mathbf{E}\) , one sets \(\mathrm{N}_p(f) = m^* (|f|^p)^{1 / p}\) and one denotes by \(\mathcal{F}^p\) the set of functions \(f\) such that \(\mathrm{N}_p(f)\) is finite (cf.~§1, Exer. 4). One equips the vector space \(\mathcal{F}^p\) with the semi-norm \(\mathrm{N}_p\) , for which it is complete.
\end{enumerate}

\begin{enumerate}
\def\labelenumi{\alph{enumi})}
\item
  Let \(\mathcal{N}\) be the set of functions f such that \(m^{*}(|f|)=0\) , \(\mathcal{V}\) the vector space generated by the finite functions belonging to \(\mathcal{M}\) , and \(\mathcal{L}^{p}\) the smallest closed linear subspace of \(\mathcal{F}^{p}\) containing the characteristic functions of the sets \(A\in \mathfrak{T}\) such that \(m^{*}(A)\) is finite. Show that \(\mathcal{L}^{p}=\mathcal{N}+(\mathcal{V}\cap\mathcal{F}^{p})\) .
\item
  Extend Lebesgue's theorem (Ch. IV, §3, No.~7, Th. 6) to the present case, and examine in particular the case \(p = 1\) (one will define the integral of a function \(f \in \mathcal{L}^1\) ).
\item
  A subset \(\mathbf{A}\) of \(\mathbf{E}\) is said to be \(m\) -integrable if \(\varphi_{\mathbf{A}} \in \mathcal{L}^1\) . Show that there then exists a set \(\mathbf{A}' \in \mathfrak{T}\) such that \(m^{*}(\mathbf{A} \cap \mathbb{C}\mathbf{A}') = m^{*}(\mathbf{A}' \cap \mathbb{C}\mathbf{A}) = 0\) . Converse, in the case that \(m(\mathbf{E})\) is finite.
\end{enumerate}

¶ 6) Let E be a set, T a tribe of subsets of E, and m an abstract measure on (E, T). Suppose that \(m(E)\) is finite and that, for every \(B \in T\) such that \(m(B) > 0\) , there exists a set \(A \in T\) such that \(A \subset B\) and \(0 < m(A) < m(B)\) . For every \(B \in T\) and every number t such that \(0 \leqslant t \leqslant m(B)\) , there exists a set \(A \in T\) such that \(m(A) = t\) and \(A \subset B\) .

¶ 7) Let T be a topological space, \(\Phi\) a clan of subsets of T, and \(m\) a relatively bounded additive mapping (cf.~Exer. 1) of \(\Phi\) into R. Assume that there exists an open covering \(\mathfrak{U}\) of T such that \(\Phi\) consists of the Borel subsets of T contained in the union of a finite number of elements of \(\mathfrak{U}\) . Assume, moreover, that for every \(A \in \Phi\) and every \(\varepsilon > 0\) , there exist a compact subset K of T and an open subset U of T such that \(K \subset A \subset U\) , \(U \in \Phi\) and \(|m|(U - K) < \varepsilon\) . Show that there exists a measure \(\mu\) on T (not necessarily positive) such that \(m(A) = \mu^{\bullet}(A)\) for all \(A \in \Phi\) , and that such a measure \(\mu\) is unique (on introducing the decomposition \(m = m_{+} - m_{-}\) of Exer. 1 c), reduce to the case that \(m\) takes positive values; first treat the case that \(T \in \Phi\) and observe that \(m\) is then inner regular; finally, treat the general case by pasting together measures).

\begin{enumerate}
\def\labelenumi{\arabic{enumi})}
\setcounter{enumi}{7}
\tightlist
\item
  Let T be a topological space and m a bounded, countably additive (cf.~Exer. 1d)) mapping of \(\mathfrak{B}(\mathrm{T})\) into R. Denote by T the set of Borel subsets A of T having the following property: for every \(\varepsilon > 0\) , there exist a closed set F and an open set U in T such that \(F \subset A \subset U\) and \(|m(B)| < \varepsilon\) for every Borel subset B of U - F (in other words, \(|m|(\mathrm{U} - \mathrm{F}) < \varepsilon\) ). Show that T is a tribe of subsets of T (first observe that \(T \in T\) and that T is a clan; it then suffices to prove that if \((\mathrm{A}_{n})_{n \in \mathbb{N}}\) is a sequence of pairwise disjoint elements of T, the set \(A = \bigcup_{n \in N} A_{n}\) belongs to T; to this end, choose closed sets \(F_{n}\) and open sets \(U_{n}\) such that \(F_{n} \subset A_{n} \subset U_{n}\) and \(|m|(\mathrm{U}_{n} - \mathrm{F}_{n}) < \varepsilon / 2^{n}\) , then set \(F = F_{0} \cup \cdots \cup F_{p}\) for p sufficiently large and \(U = \bigcup_{n \in \mathbf{N}} U_{n}\) .
\end{enumerate}


\begin{enumerate}
\def\labelenumi{\arabic{enumi})}
\setcounter{enumi}{8}
\tightlist
\item
  Let \(\mathbf{X}\) be a set; one calls gauge on \(\mathbf{X}\) every countably additive mapping of \(\mathfrak{P}(\mathbf{X})\) into \(\mathbf{R}_+\) ; a gauge \(m\) on \(\mathbf{X}\) is said to be diffuse if \(m(\{x\}) = 0\) for all \(x \in \mathbf{X}\) , and to be atomic if there exists a positive function \(f\) on \(\mathbf{X}\) such that \(m(\mathbf{A}) = \sum_{x \in \mathbf{A}} f(x)\) for every
\end{enumerate}

subset A of X.

\begin{enumerate}
\def\labelenumi{\alph{enumi})}
\item
  Show that every gauge on X may be decomposed in a unique way as the sum of a diffuse gauge and an atomic gauge.
\item
  Let \(X'\) be a subset of \(X\) , and \(m'\) a gauge on \(X'\) ; for every subset \(A\) of \(X\) , one sets \(m(A) = m'(A \cap X')\) . Show that \(m\) is a gauge on \(X\) and that it is diffuse (resp. atomic) if and only if \(m'\) has the same property.
\item
  Let \(m\) be a gauge on \(X\) and let \((X_i)_{i \in I}\) be a family of pairwise disjoint subsets of \(X\) . If every gauge on \(I\) is atomic, then \(m\big(\bigcup_{i \in I} X_i\big) = \sum_{i \in I} m(X_i)\) .
\end{enumerate}

\begin{enumerate}
\def\labelenumi{\arabic{enumi})}
\setcounter{enumi}{9}
\tightlist
\item
  Let \(\mathbf{X}\) be an uncountable infinite set, equipped with a well-ordering relation denoted \(x \leqslant y\) . For every \(x \in \mathbf{X}\) , one denotes by \(\mathrm{I}(x)\) the set of \(y \in \mathbf{X}\) such that \(y < x\) . We make the following hypotheses: \(\alpha\) ) there exists a largest element \(a\) in \(\mathbf{X}\) ; \(\beta\) ) the set of cardinals strictly smaller than \(\operatorname{Card}(\mathbf{X})\) has a largest element \(\mathfrak{c}\) ; \(\gamma\) ) for every \(x\) in \(\mathrm{I}(a)\) , one has \(\operatorname{Card}(\mathrm{I}(x)) \leqslant \mathfrak{c}\) , and the set \(\mathbf{Y}\) of \(x \in \mathbf{X}\) such that \(\operatorname{Card}(\mathrm{I}(x)) = \mathfrak{c}\) is nonempty;
\end{enumerate}

\(\delta)\) every gauge on a set with cardinal \(\leqslant \mathfrak{c}\) is atomic. Show that every gauge on X is atomic. One can argue by contradiction in the following way: let \(m\) be a nonzero diffuse gauge on X. Denote by \(b\) the smallest element of Y, and for every \(x \in I(a)\) let \(f_x\) be an injection of \(I(x)\) into \(I(b)\) . For every pair \((x,y) \in I(a) \times I(b)\) , denote by \(A_{x,y}\) the set of elements \(z \in X\) such that \(x < z < a\) and \(f_z(x) = y\) ; for every \(x \in I(a)\) , denote by \(M_x\) the set of \(y \in I(b)\) such that \(m(A_{x,y}) > 0\) , and for every \(y \in I(b)\) , denote by \(N_y\) the set of \(x \in I(a)\) such that \(m(A_{x,y}) > 0\) . Using Exer. 9 c), show that \(m(X) = \sum_{y \in I(b)} m(A_{x,y})\) ; from this, deduce that the set \(M_x\) is countable and nonempty for every \(x \in I(a)\) ; then show that each of the sets \(N_y\) is countable and from this deduce the contradiction \(\text{Card}(X) = c\) .

\begin{enumerate}
\def\labelenumi{\arabic{enumi})}
\setcounter{enumi}{10}
\tightlist
\item
  A cardinal \(\mathfrak{c}\) is said to be an Ulam cardinal (or to be ulamian) if every gauge on a set with cardinal \(\mathfrak{c}\) is atomic. Show that every countable cardinal is ulamian, that if \(\mathfrak{c}\) is an Ulam cardinal then so is every cardinal \(\mathfrak{c}' \leqslant \mathfrak{c}\) , and that if \((\mathfrak{c}_i)_{i \in I}\) is a family of Ulam cardinals such that \(\operatorname{Card}(I)\) is ulamian, then the cardinal \(\sum_{i \in I} \mathfrak{c}_i\) is ulamian. Finally,
\end{enumerate}

the smallest uncountable cardinal is ulamian (cf.~Exer. 10).

\begin{enumerate}
\def\labelenumi{\arabic{enumi})}
\setcounter{enumi}{11}
\tightlist
\item
  Let X be a set, U a subset of \(\mathfrak{P}(X)\) , and m the characteristic function of U. For m to be a gauge, it is necessary and sufficient that U be an ultrafilter and that the intersection of every countable subset of U belong to U. Under these conditions, U is said to be an Ulam ultrafilter on X. Suppose that this is the case, and that \((\mathrm{H}_{i})_{i\in\mathrm{I}}\) is a family of elements of U such that \(\operatorname{Card}(\mathrm{I})\) is an Ulam cardinal; show that \(\bigcap_{i\in I} H_{i}\) belongs
\end{enumerate}

to \(\mathfrak{U}\) .

\begin{enumerate}
\def\labelenumi{\arabic{enumi})}
\setcounter{enumi}{12}
\tightlist
\item
  In this exercise, we admit the continuum hypothesis, that is, we adjoin to the axioms of set theory the axiom \(\langle\mathrm{Card}(\mathbf{R})\) is the smallest uncountable cardinal \(\rangle\) . \(^{(1)}\)
\end{enumerate}

\begin{enumerate}
\def\labelenumi{\alph{enumi})}
\item
  Every gauge on \(\mathbf{R}\) is atomic (cf.~Exer. 11).
\item
  Let X be a set and m a gauge on X satisfying the following property: for every subset A of X such that \(m(A) > 0\) , there exists a subset B of A such that \(0 < m(B) < m(A)\) ; then m is zero (for every integer \(n \geqslant 1\) , there exists a finite partition \((\mathrm{X}_{n,i})_{i \in \mathrm{I}_n}\) of X such that \(m(\mathrm{X}_{n,i}) \leqslant 1/n\) for all \(i \in I_n\) ; set \(I = \prod_{n \geqslant 1} I_n\) and
\end{enumerate}

\(\mathrm{X}_{\alpha} = \bigcap_{n\geqslant 1}\mathrm{X}_{n,\alpha_n}\) for every \(\alpha = (\alpha_{n})_{n\geqslant 1}\) in I; the family \((\mathrm{X}_i)_{i\in \mathrm{I}}\) is a partition of \(\mathbf{X}\) ,

\(m(X_{i}) = 0\) for every \(i \in I\) , and every gauge on \(I\) is atomic by \(a\) ); conclude by means of Exer. 9 c)).

\begin{enumerate}
\def\labelenumi{\alph{enumi})}
\setcounter{enumi}{2}
\item
  Let X be a set on which there does not exist any nontrivial Ulam ultrafilter; then every gauge on X is atomic (otherwise, by b), there would exist a diffuse gauge m on X and a subset A of X such that \(m(A) > 0\) and such that \(m(B) = 0\) or \(m(A - B) = 0\) for every subset B of A; the set U of subsets B of X such that \(m(A) = m(A \cap B)\) would be a nontrivial Ulam ultrafilter on X).
\item
  Show that \(2^{\mathfrak{c}}\) is an Ulam cardinal for every Ulam cardinal \(\mathfrak{c}\) (by \(c\) ), it suffices to show that if \(C\) is a set with cardinal \(\mathfrak{c}\) , then every Ulam ultrafilter \(\mathfrak{U}\) on \(X = \{0,1\}^{C}\) is trivial; equip \(C\) with a well-ordering structure and define by transfinite induction a family \(s = (s_{\alpha})_{\alpha \in C}\) of elements of \(\{0,1\}\) such that set of \(x = (x_{\alpha})_{\alpha \in C}\) in \(X\) for which \(x_{\alpha}=s_{\alpha}\) for all \(\alpha\leqslant\beta\) belongs to \(\mathfrak{U}\) for every \(\beta\in\mathrm{C}\) (cf.~Exer. 12); then \(s\) belongs to every element of \(\mathfrak{U}\) ). \(^{(2)}\)
\end{enumerate}

\begin{enumerate}
\def\labelenumi{\arabic{enumi})}
\setcounter{enumi}{13}
\item
  Let \((\mathbf{T}_i)_{i\in I}\) be a family of Radon topological spaces, and \(\mathbf{T}\) the sum space of the family. If \(\operatorname{Card}(\mathbf{I})\) is an Ulam cardinal, then the topological space \(\mathbf{T}\) is a Radon space (show that for every positive and bounded countably additive function \(m\) on the Borel tribe of \(\mathbf{T}\) , there exists a countable subset \(J\) of \(I\) such that \(m(\bigcup_{i \in I - J} \mathbf{T}_i) = 0\) .
\item
  Let T be a discrete space whose cardinal is the smallest uncountable cardinal. Show that T is a Radon locally compact space whose topology does not have a countable base (cf.~Exer. 11). Deduce from this that there exist compact Radon spaces that are not metrizable.
\item
  Let \(I\) be a countable set, and \((\mathbf{T}_i)_{i \in I}\) a family of Radon spaces. Assume that every compact subset of any of the spaces \(\mathbf{T}_i\) is metrizable. Show that the product space \(\mathbf{T} = \prod_{i \in I} \mathbf{T}_i\) is Radon. Generalize to the case of inverse limits.
\item
  Let \(\mathbf{T}\) be a topological space.
\end{enumerate}

\begin{enumerate}
\def\labelenumi{\alph{enumi})}
\tightlist
\item
  Let \(m\) be a countably additive mapping of \(\mathfrak{B}(\mathbf{T})\) into \(\mathbf{R}_+\) . Assume that there exists a sequence \((\mathbf{T}_n)_{n\in \mathbf{N}}\) of Borel subsets of \(\mathbf{T}\) such that \(\mathbf{T} = \bigcup_{n\in \mathbf{N}}\mathbf{T}_n\) , \(m(\mathbf{T}_0) = 0\)
\end{enumerate}

and the subspace \(\mathbf{T}_n\) of \(\mathbf{T}\) is Radon for every \(n \geqslant 1\) . Show that there exists a bounded positive measure \(\mu\) on \(\mathbf{T}\) such that \(m(\mathbf{A}) = \mu^{\bullet}(\mathbf{A})\) for every Borel subset \(\mathbf{A}\) of \(\mathbf{T}\) (one can reduce to the case that the \(\mathbf{T}_n\) are pairwise disjoint on applying the Cor. of Prop. 2 of No.~3).

\begin{enumerate}
\def\labelenumi{\alph{enumi})}
\setcounter{enumi}{1}
\item
  Extend the preceding result to the case that the sets \(T_{n}\) for \(n \geqslant 1\) are no longer assumed Borel but merely universally measurable in T (argue as in Prop. 2 of No.~3).
\item
  Every topological space that is the union of a sequence of metrizable compact subspaces is Radon.
\end{enumerate}

\begin{enumerate}
\def\labelenumi{\arabic{enumi})}
\setcounter{enumi}{17}
\tightlist
\item
  Let \(\mathbf{T}_1\) and \(\mathbf{T}_2\) be two topological spaces and \(f\) a bijective continuous mapping of \(\mathbf{T}_1\) onto \(\mathbf{T}_2\) . Assume that, for every bounded countably additive function \(m_1\) on the Borel tribe \(\mathfrak{B}(\mathbf{T}_1)\) , there exists a sequence of compact subsets \(\mathrm{K}_n\) of \(\mathbf{T}_1\) such that \(m_1(\mathbf{T}_1 - \bigcup_{n}\mathbf{K}_n) = 0\) . Let \(m_2\) be a bounded countably additive function on \(\mathfrak{B}(\mathbf{T}_2)\) . In
\end{enumerate}

order that there exist a bounded countably additive function \(m_{1}\) on \(\mathfrak{B}(\mathrm{T}_{1})\) such that \(m_{2}(\mathrm{A}) = m_{1}(f^{-1}(\mathrm{A}))\) for all \(\mathrm{A} \in \mathfrak{B}(\mathrm{T}_{2})\) , it is necessary and sufficient that the following condition be fulfilled: there exists a sequence of compact subsets \(K_{n}\) of \(T_{1}\) such that \(m_{2}(\mathrm{T}_{2} - \bigcup_{n} f(\mathrm{K}_{n})) = 0\) .

